\documentclass[11pt,a4paper]{article}

\usepackage[utf8]{inputenc}
\usepackage[T1]{fontenc}
\usepackage[spanish,es-noindentfirst]{babel}
\usepackage[margin=2.2cm]{geometry}
\usepackage{xcolor}
\definecolor{brandblue}{HTML}{0B3C6B}
\definecolor{brandgray}{HTML}{4A4A4A}
\definecolor{tablehead}{HTML}{E8EEF4}
\usepackage{hyperref}
\hypersetup{
  colorlinks=true,
  linkcolor=brandblue,
  urlcolor=brandblue,
  citecolor=brandblue,
  pdftitle={Analisis tecnico Resility / Zentris},
  pdfauthor={Joaquin Cuevas}
}
\usepackage{booktabs}
\usepackage{longtable}
\usepackage{array}
\usepackage{enumitem}
\usepackage{titlesec}
\usepackage{fancyhdr}
\usepackage{parskip}
\usepackage{graphicx}

\titleformat{\section}{\large\bfseries\color{brandblue}}{\thesection.}{0.5em}{}
\titleformat{\subsection}{\bfseries\color{brandblue}}{\thesubsection}{0.5em}{}

\pagestyle{fancy}
\fancyhf{}
\renewcommand{\headrulewidth}{0.4pt}
\fancyhead[L]{\small\color{brandgray} Resility / Zentris -- Analisis tecnico}
\fancyhead[R]{\small\color{brandgray} \thepage}
\fancyfoot[C]{\small\color{brandgray} Informe de uso interno -- OSINT pasivo, 2026-11-20}

\newcolumntype{L}[1]{>{\raggedright\arraybackslash}p{#1}}

\title{\color{brandblue}\textbf{Resility / Zentris}\\[0.3em]
\Large Analisis tecnico de seguridad e inteligencia artificial\\[0.2em]
\large Informe de referencia para hackathon de Ciberseguridad + IA (SOC / agentes)}
\author{Joaquin Cuevas}
\date{20 de noviembre de 2026}

\begin{document}
\maketitle
\thispagestyle{fancy}

\begin{center}
\colorbox{tablehead}{\parbox{0.92\textwidth}{\small
\textbf{Nota de alcance:} este informe se construyo exclusivamente con reconocimiento
pasivo sobre informacion publica de \url{resility.io} (HTML, cabeceras HTTP, \texttt{robots.txt},
datos estructurados JSON-LD, JavaScript servido publicamente, y descripciones publicas del
producto). No se realizaron pruebas activas, escaneos de vulnerabilidades ni ningun tipo de
interaccion con la plataforma Zentris en si, que es un producto SaaS para clientes y no una
aplicacion descargable. Las cifras de desempeno citadas son afirmaciones del proveedor, no
verificadas de forma independiente.}}
\end{center}

\vspace{0.5em}
\tableofcontents
\vspace{1em}

% ---------------------------------------------------------------------------
\section{Resumen ejecutivo}

\textbf{Resility} es una empresa chilena de ciberseguridad gestionada (MSSP), fundada en 2016,
con SOC propio 24/7 en Santiago y presencia en Chile, Peru, Colombia y Espana. Su diferenciador
competitivo es \textbf{Zentris}, una plataforma propia de agentes de IA (``SOC agentico'') que
automatiza triage de alertas, gestion continua de exposicion y pentesting continuo sobre el
stack de seguridad que el cliente ya tiene, sin reemplazar EDR/SIEM existentes (modelo
\emph{agentless}).

Para una hackathon de \textbf{Ciberseguridad + IA / agentes de SOC}, Zentris es un caso de
referencia util en tres sentidos:

\begin{enumerate}[leftmargin=1.4em]
  \item Es un \textbf{patron de arquitectura validado comercialmente}: pipeline de ingesta
  $\rightarrow$ triage multiagente $\rightarrow$ validacion humana (\emph{Expert-in-the-loop},
  EITL) $\rightarrow$ respuesta. Sirve como blueprint de partida.
  \item Expone \textbf{vacios publicos concretos} (gobernanza de agentes ofensivos autonomos,
  explicabilidad de decisiones, aislamiento multi-tenant, metodologia de metricas) que son
  exactamente el tipo de problema que un equipo de hackathon puede resolver y demostrar en
  una demo corta.
  \item Su propio sitio muestra \textbf{buenas practicas de higiene tecnica} (CSP estricta,
  HSTS preload, \texttt{robots.txt} con \emph{content signals} anti-entrenamiento de IA) que
  sirven de checklist minimo para cualquier proyecto que se presente.
\end{enumerate}

% ---------------------------------------------------------------------------
\section{Metodologia}

\begin{itemize}[leftmargin=1.4em]
  \item \textbf{Fuente de verdad:} unicamente contenido publico (sitio web, \texttt{robots.txt},
  JSON-LD, JavaScript estatico, FAQ, material de prensa/ProChile referenciado por el propio sitio).
  \item \textbf{Herramienta evaluada:} \href{https://github.com/morluto/rea}{REA (``Reverse
  Engineer Anything'')}, una utilidad MIT de proposito general para ingenieria inversa de
  binarios, apps Electron/JS y sitios web. No contenia nada especifico de Resility; se usa aqui
  solo como referencia metodologica (su guia de ``browser observation'' para sitios web),
  aplicada manualmente.
  \item \textbf{Sin pruebas activas:} no se ejecutaron escaneos de puertos, fuzzing, intentos de
  autenticacion, ni se interactuo con la plataforma Zentris (no es accesible publicamente).
  \item \textbf{Evidencia cruda} disponible en el repositorio: cabeceras HTTP, \texttt{robots.txt},
  JSON-LD y una muestra de requests de red observados al cargar la pagina de inicio.
\end{itemize}

% ---------------------------------------------------------------------------
\section{Resility: panorama de la empresa}

\begin{longtable}{L{4cm}L{11cm}}
\toprule
\textbf{Dato} & \textbf{Detalle} \\
\midrule
\endhead
Fundacion & 2016 (``10 anos de servicio'') \\
Sede / SOC & Santiago, Chile (SOC propio, operacion 24x7x365) \\
Presencia & Chile, Peru, Colombia, Espana \\
Certificacion & ISO/IEC 27001:2022 (SGSI certificado, auditado) \\
Marco regulatorio objetivo & Ley 21.663 (Marco de Ciberseguridad / ANCI, reporte en 3h para
Operadores de Importancia Vital), Ley 21.719 (datos personales), normativa CMF (banca, fintech,
seguros, cooperativas) \\
Equipo & 14 especialistas en operacion/ofensiva/desarrollo; 71 certificaciones del equipo;
niveles N1, N2 y EITL (Expert-in-the-loop) \\
Grupo & 3 empresas bajo un mismo ecosistema: Resility (servicios gestionados), Zentris (producto
de IA propio), Staff21x (agentes de IA para automatizacion B2B general) \\
Clientes citados & Sector financiero, mineria, salud, retail e industria critica (logos
publicados: Gasco, Collahuasi, CAP, Red Salud, Autofin, Capual, Talana, entre otros) \\
Reconocimiento & Zentris seleccionado por ProChile para misiones comerciales (Colombia, Peru,
Nueva York, Portugal) \\
\bottomrule
\end{longtable}

\subsection{Servicios ofrecidos (lineas de negocio)}

\begin{longtable}{L{4cm}L{11cm}}
\toprule
\textbf{Linea} & \textbf{Servicios} \\
\midrule
\endhead
Gestionados 24x7 & MDR gestionado, CyberSOC 24x7 o 5x8, NOC-as-a-Service, gestion de
vulnerabilidades (Tenable), CAASM (inventario de ciberactivos), gestion de endpoints/EDR,
gestion de firewalls perimetrales \\
Profesionales / ofensivos & Ethical hacking, pentesting de apps web/moviles/APIs, red/purple
team, re-testing tecnico, respuesta a incidentes (IR), analisis forense digital, salas de crisis
/ TTX, ethical phishing \\
Consultoria y cumplimiento & Preparacion ANCI (Ley 21.663), adecuacion Ley 21.719, implementacion
y auditoria ISO 27001, diagnostico CIS v8, normativa CMF \\
Capacitacion & Resility Academy (SOC N1, threat hunting, IR forense, cloud security), programas
de \emph{awareness} (phishing simulado, ingenieria social, programa C-level) \\
\bottomrule
\end{longtable}

\subsection{Stack tecnologico declarado}

Wazuh (SIEM/XDR open source), Tenable/Nessus (exposure management), Bitdefender GravityZone
(EDR), Grafana Enterprise (dashboards), Akamai (nube de alta disponibilidad, 99.9\% uptime
reclamado para su plataforma), Fortinet / Check Point / Palo Alto / pfSense (firewalls),
Microsoft Sentinel / Microsoft 365, y MITRE ATT\&CK\textregistered{} como marco de mapeo de
detecciones.

% ---------------------------------------------------------------------------
\section{Zentris: arquitectura de IA (segun fuentes publicas)}

Zentris se presenta como ``SOC agentico'': agentes de IA que investigan, clasifican y responden
sobre el EDR/SIEM existente del cliente, mas un modulo de pentesting continuo. Es una plataforma
\textbf{agentless} (se conecta por API en la nube, sin instalar agentes pesados ni hardware).

\subsection{Pipeline de operacion}

\begin{longtable}{L{2.6cm}L{12.4cm}}
\toprule
\textbf{Etapa} & \textbf{Descripcion publica} \\
\midrule
\endhead
1. Recolecta & Correlacion de eventos SIEM/XDR: endpoints EDR, firewalls, nube AWS/Azure/GCP,
Active Directory, Microsoft 365. \\
2. Investiga & ``Agentes Zentris'': triage autonomo y enriquecimiento de alertas con inteligencia
de amenazas ``a velocidad de maquina''. \\
3. Valida & Analistas humanos N1/N2 + \textbf{EITL} (Expert-in-the-loop) para descartar falsos
positivos y evaluar contexto critico. \\
4. Responde & Contencion activa (aislamiento de hosts comprometidos), escalamiento inmediato,
reporte formal segun SLA. \\
\bottomrule
\end{longtable}

\subsection{Modulos del producto}

\begin{itemize}[leftmargin=1.4em]
  \item \textbf{Pentesting continuo 24x7}: ``agentes ofensivos sin pausa''; al corregir una
  vulnerabilidad la alerta se cierra automaticamente (sin re-test manual programado).
  \item \textbf{Exposure management}: revision continua de sitios y superficie expuesta a
  internet, declarado ``sin tope de URLs''.
  \item \textbf{SOC agentico}: investigacion/clasificacion/respuesta automatizada sobre el
  EDR/SIEM existente.
  \item \textbf{NOC unificado}: disponibilidad y metricas de desempeno en el mismo panel.
  \item \textbf{Cumplimiento y reporteria en espanol}: generacion automatica de evidencia para
  Ley 21.663 (ANCI), ISO 27001, CMF y NIST CSF, con vision ejecutiva para directorio/regulador.
  \item \textbf{Integraciones de alertamiento}: Slack, Microsoft Teams, Discord, Webhooks.
  \item \textbf{Despliegue agentless}: 24--48 horas, via API en la nube.
\end{itemize}

\subsection{Metricas de marketing (no auditadas)}

\begin{longtable}{L{6cm}L{8cm}}
\toprule
\textbf{Metrica} & \textbf{Valor declarado} \\
\midrule
\endhead
Tiempo de deteccion & \texttt{< 15 min} \\
Reduccion de falsos positivos & \texttt{-93\%} \\
Tiempo de despliegue & \texttt{24--48 h} \\
Costo vs. SOC tradicional & \texttt{-85\%} \\
Composicion de agentes & ``\textasciitilde 20 agentes especializados autonomos'' \\
\bottomrule
\end{longtable}

\textbf{Importante:} no hay metodologia publica (baseline, tamano de muestra, si son promedios
de clientes reales) detras de estas cifras. Tratarlas como afirmaciones de marketing, no como
benchmarks verificados.

% ---------------------------------------------------------------------------
\section{Postura de seguridad del sitio publico}

Evidencia cruda en el repositorio: \texttt{evidence/headers-resility-io.txt},
\texttt{evidence/robots.txt}, \texttt{evidence/structured-data.json},
\texttt{evidence/network-requests-sample.txt}.

\begin{longtable}{L{4cm}L{11cm}}
\toprule
\textbf{Aspecto} & \textbf{Hallazgo} \\
\midrule
\endhead
Hosting / CDN & Cloudflare (\texttt{Server: cloudflare}, beacon RUM en \texttt{/cdn-cgi/rum}).
Sitio generado con \textbf{Astro} (\texttt{/\_astro/*.css}), con al menos una ruta de API propia
(\texttt{/api/contact}, detectada en \texttt{site.js}) -- probablemente Cloudflare
Pages/Workers/Functions. \\
Cabeceras de seguridad & HSTS con \texttt{preload} e \texttt{includeSubDomains};
\texttt{Content-Security-Policy} estricta (\texttt{default-src 'self'}, sin
\texttt{unsafe-inline} en scripts, usa hash \texttt{sha256-...} para el unico inline permitido);
\texttt{X-Frame-Options: SAMEORIGIN}; \texttt{X-Content-Type-Options: nosniff};
\texttt{Cross-Origin-Opener-Policy: same-origin}; \texttt{Permissions-Policy} restrictiva
(geolocalizacion, camara, microfono, pagos y USB deshabilitados). \\
Anti-bot & Cloudflare Turnstile + captcha matematico simple en el formulario de contacto
(doble capa). \\
Certificado & Sectigo PositiveSSL. \\
Tracking & Meta Pixel, Google Ads, LinkedIn Insight Tag, Google Tag Manager -- declarados
explicitamente como dominios permitidos en la CSP. \\
\texttt{robots.txt} & Usa \emph{Content Signals} (\texttt{search=yes, ai-train=no,
use=reference}) y bloquea explicitamente crawlers de entrenamiento de IA (GPTBot, ClaudeBot,
Google-Extended, Bytespider, CCBot, etc.), permitiendo indexacion de busqueda normal. \\
Datos estructurados & JSON-LD declara \texttt{Organization}, \texttt{WebSite}, \texttt{FAQPage}
y un \texttt{SoftwareApplication} para Zentris (\texttt{applicationCategory:
SecurityApplication}, \texttt{operatingSystem: Cloud / Agentless}). \\
\bottomrule
\end{longtable}

\textbf{Lectura:} la higiene del sitio de marketing es solida y coherente con el discurso de
``predicar con el ejemplo''. Pero esto certifica solo el sitio publico, no dice nada sobre la
seguridad real del backend de Zentris (multi-tenencia, aislamiento entre clientes, gestion de
credenciales de integracion con el SIEM/EDR de cada cliente), que no es observable desde afuera.

% ---------------------------------------------------------------------------
\section{Vacios y oportunidades identificadas}

\begin{enumerate}[leftmargin=1.4em]
  \item \textbf{Gobernanza de agentes ofensivos autonomos}: no hay informacion publica sobre
  guardrails, sandboxing o reglas de engagement para los ``agentes de pentesting continuo''
  que actuan 24/7 sobre produccion del cliente. Riesgo de disenio: falsos positivos destructivos,
  necesidad de ventanas/alcance controlado.
  \item \textbf{Explicabilidad}: no se publica como los agentes de triage justifican una
  decision (evidencia, cadena de razonamiento, trazabilidad por regulacion ANCI/CMF).
  \item \textbf{Orquestacion multiagente}: ``\textasciitilde 20 agentes'' sin detalle de que
  framework de orquestacion, que modelo(s) LLM, ni como se coordinan o resuelven conflictos
  entre agentes.
  \item \textbf{Aislamiento multi-tenant}: sin informacion publica sobre como se separan datos
  y credenciales de distintos clientes dentro de una plataforma agentless compartida.
  \item \textbf{Metodologia de metricas}: SLA, \% de falsos positivos y ahorro de costos son
  afirmaciones sin benchmark reproducible.
\end{enumerate}

% ---------------------------------------------------------------------------
\section{\color{brandblue} Para la hackathon: Ciberseguridad + IA (SOC / agentes)}

Esta seccion traduce los hallazgos anteriores en material directamente utilizable para una
hackathon centrada en agentes de IA para SOC/ciberseguridad.

\subsection{Blueprint de referencia (derivado de Zentris)}

Un pipeline minimo viable para una demo de ``SOC agentico'', inspirado en el patron observado:

\begin{enumerate}[leftmargin=1.4em]
  \item \textbf{Ingesta}: conector agentless hacia una fuente de eventos (puede ser un SIEM
  real tipo Wazuh/OSSEC, o un dataset simulado de alertas para la demo).
  \item \textbf{Triage multiagente}: 2--4 agentes especializados (ej. agente de
  enriquecimiento de IOC, agente de correlacion MITRE ATT\&CK, agente de priorizacion por
  riesgo) que procesan la alerta en paralelo.
  \item \textbf{Capa de evidencia}: cada agente debe producir una salida con \emph{evidencia
  citada} (que regla disparo, que indicador, que confianza) -- esto es justamente lo que
  Zentris \textbf{no} publica y es un diferenciador facil de demostrar en una demo.
  \item \textbf{Human-in-the-loop (EITL)}: un paso de aprobacion/veto humano simulado antes de
  cualquier accion de contencion, con la evidencia de la etapa anterior como input.
  \item \textbf{Respuesta}: accion simulada (aislar host, abrir ticket, notificar por
  Slack/Webhook) con log auditable.
\end{enumerate}

\subsection{Checklist tecnico para evaluar/presentar un proyecto de ``agente de SOC''}

\begin{longtable}{L{4.5cm}L{10.5cm}}
\toprule
\textbf{Criterio} & \textbf{Por que importa (y que casi nadie publica)} \\
\midrule
\endhead
Guardrails de acciones autonomas & Limites explicitos de lo que un agente puede hacer sin
aprobacion humana (especialmente en pentesting u ofensiva autonoma). \\
Explicabilidad / evidencia & Cada decision deberia venir con la evidencia que la sustenta, no
solo un veredicto. \\
Reduccion de falsos positivos medible & Definir baseline y metodologia antes de prometer un
\%. \\
Human-in-the-loop real & Punto de control humano para acciones criticas o irreversibles. \\
Trazabilidad para cumplimiento & Logs/evidencia exportable mapeada a un marco (MITRE ATT\&CK,
ISO 27001, NIST CSF) -- vende bien a jurados con perfil empresarial. \\
Aislamiento / multi-tenencia & Si el proyecto es multi-cliente, mostrar como se separan datos
y credenciales. \\
Costo de integracion & Favorecer modelos ``agentless'' (API) sobre instalacion de agentes
pesados -- reduce friccion de adopcion, tal como vende Zentris. \\
\bottomrule
\end{longtable}

\subsection{Ideas de proyecto concretas}

\begin{itemize}[leftmargin=1.4em]
  \item \textbf{``Evidencia primero''}: un agente de triage que, a diferencia de lo que Zentris
  publica, muestre explicitamente su cadena de evidencia (regla MITRE ATT\&CK, IOC, nivel de
  confianza) para cada alerta clasificada -- resuelve el vacio de explicabilidad identificado
  en la seccion 6.
  \item \textbf{Guardrails para pentesting autonomo}: un framework ligero de ``reglas de
  engagement'' (RoE) programables -- ventanas horarias, alcance de IPs/dominios, severidad
  maxima de explotacion permitida -- que envuelva cualquier agente ofensivo autonomo antes de
  ejecutar una accion contra un objetivo real.
  \item \textbf{Validador de metricas de SOC con IA}: una herramienta que, dado un log de
  alertas reales vs. clasificacion de un agente, calcule tasa de falsos positivos/negativos de
  forma reproducible -- justo la metodologia que falta en el marketing de la industria (y de
  Zentris).
  \item \textbf{Reporte de cumplimiento automatico}: generador que tome eventos de seguridad y
  los mapee automaticamente a evidencia para Ley 21.663 (ANCI) / ISO 27001 -- demo rapida y con
  gancho regulatorio local (Chile).
  \item \textbf{Mini-SOC multiagente open source}: usando Wazuh (gratis) + un orquestador de
  agentes (framework de agentes LLM a eleccion del equipo) + un dashboard tipo Grafana, para
  reproducir el patron de Zentris con stack 100\% abierto.
\end{itemize}

\subsection{Consideraciones eticas y de seguridad al construir esto}

\begin{itemize}[leftmargin=1.4em]
  \item Un agente de ``pentesting autonomo'' nunca deberia ejecutarse contra infraestructura
  real sin autorizacion explicita y alcance acotado (reglas de engagement) -- demostrarlo contra
  un entorno de laboratorio propio (ej. una VM vulnerable a proposito) es suficiente para una
  demo y evita riesgos legales.
  \item Si el proyecto usa LLMs para decisiones de seguridad, documentar que el modelo
  \emph{no} ejecuta acciones irreversibles sin aprobacion humana (patron EITL) -- es lo que el
  propio mercado (incluido Zentris) ya reconoce como necesario.
  \item Evitar prometer metricas de reduccion de falsos positivos sin mostrar la metodologia:
  es el error de comunicacion mas comun en este espacio y fue justamente uno de los vacios
  identificados en el analisis.
\end{itemize}

% ---------------------------------------------------------------------------
\section{Conclusion}

Resility es un MSSP chileno establecido que diferencia su oferta con Zentris, una capa de
agentes de IA para triage, exposure management y pentesting continuo sobre el stack de
seguridad existente del cliente (modelo agentless). La propuesta es coherente con la tendencia
de ``SOC aumentado por IA con humano en el loop'' y esta respaldada por integraciones reales
(Wazuh, Tenable, Bitdefender) y reconocimiento publico (ProChile). Sus afirmaciones de
rendimiento y el detalle interno de los agentes no son publicos ni verificables desde afuera.
Para efectos de una hackathon de ciberseguridad e IA, el valor principal de este analisis esta
en el \textbf{blueprint de arquitectura} (seccion 7.1) y en los \textbf{vacios identificados}
(seccion 6), que senalan exactamente donde un equipo puede diferenciarse con una demo concreta:
evidencia/explicabilidad, guardrails de agentes ofensivos y metodologia de metricas reproducible.

% ---------------------------------------------------------------------------
\section{Evidencia y fuentes}

\begin{itemize}[leftmargin=1.4em]
  \item Repositorio con evidencia cruda:
  \url{https://github.com/Joaquincuevas/resility-zentris-analysis}
  \item Sitio analizado: \url{https://resility.io/}
  \item Herramienta de referencia evaluada: \url{https://github.com/morluto/rea}
\end{itemize}

\vspace{1em}
\noindent\small\color{brandgray}
Informe elaborado por Joaquin Cuevas a partir de informacion publica, 20 de noviembre de 2026.
Uso previsto: material de referencia para hackathon de ciberseguridad e IA. No constituye una
auditoria de seguridad de Resility/Zentris ni una certificacion de sus controles internos.

\end{document}
